% ============================================================
% Come si usano, argomento 6: pmf condizionate e funzioni di una variabile
% ============================================================
\argomento{Pmf condizionate e funzioni di una variabile}{Formule a p.~\pageref*{f:pmfcond}}[Pmf condizionate e funzioni]

\begin{scelta}{Come scegliere la formula}
\renewcommand{\arraystretch}{1.25}
\begin{tabular}{@{}P{0.56\linewidth}P{0.38\linewidth}@{}}
  \toprule
  \textbf{Se il testo\ldots} & \textbf{si usa\ldots} \\
  \midrule
  dà un'informazione su $X$ o sul risultato («sapendo che $X > 4$»)
    & pmf e media condizionate \vv{pmfcond} \vv{mediacond} \\
  fa dipendere $X$ da un caso che non si conosce (quale moneta, quale canale)
    & probabilità totale per le pmf \vv{totalepmf}, media pesata \vv{mediamedie} \\
  chiede il caso sapendo il valore di $X$ («uscite 9 teste: quale moneta?»)
    & Bayes \vv{totalepmf}, \vv{bayes} \\
  chiede la pmf di $Y = g(X)$
    & la somma sulla controimmagine \vv{pmffunzione} \\
  chiede solo la media di $Y = g(X)$
    & il teorema fondamentale \vv{valoremedio}: la pmf di $Y$ non serve \\
  \bottomrule
\end{tabular}
\end{scelta}

\scheda{pmfcond}{
  \usoquando{Si sa che si è verificato un evento $A$, spesso definito tramite $X$ stessa («sapendo che $X > 4$»), e
    si vuole come si distribuisce $X$ con questa informazione.}
  \usocome{Si tengono solo i valori compatibili con $A$ e si dividono le loro probabilità per $P(A)$: si «rialzano»
    in proporzione, così sommano di nuovo a 1. Gli altri valori hanno probabilità 0.}
  \usoesempio{Dado onesto, sapendo che $X > 4$: $p_{X \mid X > 4}(5) = p_{X \mid X > 4}(6) =
    \frac{1/6}{2/6} = \frac{1}{2}$, e 0 per le altre facce.}
}

\scheda{totalepmf}{
  \usoquando{$X$ dipende da un caso che non si conosce, ma sapendo il caso la sua pmf è nota: per esempio una
    moneta presa a caso fra due, poi lanciata più volte.}
  \usocome{La pmf di $X$ è la media pesata delle pmf condizionate, con pesi le probabilità dei casi.}
  \usovarianti{Al contrario, osservato il valore $X = x$, la probabilità di ciascun caso si trova con Bayes: al
    numeratore il termine del caso, al denominatore $p_X(x)$.}
  \usoesempio{Una moneta onesta e una con $p = 0{,}8$, se ne prende una a caso e la si lancia 10 volte:
    $p_X(9) = \frac{1}{2} \cdot 0{,}0098 + \frac{1}{2} \cdot 0{,}2684$. Se sono uscite 9 teste, la moneta è quella
    truccata con probabilità $\frac{0{,}2684}{0{,}0098 + 0{,}2684} \simeq 0{,}965$.}
}

\scheda{mediacond}{
  \usoquando{Si vuole la media di $X$ sapendo che si è verificato l'evento $E$.}
  \usocome{Si trova la pmf condizionata \vv{pmfcond} e se ne calcola la media.}
  \usoesempio{Dado onesto: $E[X \mid X > 4] = 5 \cdot \frac{1}{2} + 6 \cdot \frac{1}{2} = 5{,}5$.}
}

\scheda{mediamedie}{
  \usoquando{La media di $X$ è difficile da calcolare direttamente, ma in ciascun caso è facile (per esempio perché,
    noto il caso, $X$ è binomiale e la sua media è $N p$).}
  \usocome{Si pesa la media di ogni caso con la probabilità del caso e si somma. Non serve la pmf di $X$.}
  \usoesempio{Le due monete dell'esempio precedente: $E[X] = \frac{1}{2} \cdot (10 \cdot 0{,}5) + \frac{1}{2} \cdot
    (10 \cdot 0{,}8) = 6{,}5$.}
}

\scheda{pmffunzione}{
  \usoquando{Si vuole la pmf di una variabile ottenuta trasformando $X$: un quadrato, una potenza $X^2 / R$, un
    seno, una quantizzazione.}
  \usocome{(1) Per ogni valore $x$ si calcola $g(x)$: i risultati, presi una volta sola, formano l'alfabeto di $Y$.
    (2) Per ogni $y$ si sommano le probabilità di tutti gli $x$ che finiscono in $y$.}
  \usoattenzione{Se $g$ non è iniettiva più valori di $X$ finiscono nello stesso $y$, e le loro probabilità vanno
    sommate: con $Y = X^2$, i valori $-2$ e $2$ danno lo stesso $y = 4$.}
  \usoesempio{$X$ uniforme su $\{-2, -1, 0, 2\}$ e $Y = X^2$: $p_Y(0) = \frac{1}{4}$, $p_Y(1) = \frac{1}{4}$,
    $p_Y(4) = \frac{1}{4} + \frac{1}{4} = \frac{1}{2}$.}
}

\scheda{valoremedio}{
  \usoquando{Si vuole solo la media di $Y = g(X)$, o un momento come $E[X^2]$: la pmf di $Y$ non serve.}
  \usocome{Si pesa ogni $g(x)$ con la probabilità $p_X(x)$ e si somma.}
  \usoattenzione{In generale $E[g(X)] \neq g(E[X])$. Se $X$ vale $\pm 1$ con probabilità $\frac{1}{2}$, allora
    $(E[X])^2 = 0$ ma $E[X^2] = 1$. La funzione si porta fuori dalla media solo se è lineare.}
  \usoesempio{$X$ uniforme su $\{-2, -1, 0, 2\}$: $E[X^2] = \frac{4 + 1 + 0 + 4}{4} = \frac{9}{4}$, lo stesso valore
    che si ottiene dalla pmf di $Y = X^2$: $0 \cdot \frac{1}{4} + 1 \cdot \frac{1}{4} + 4 \cdot \frac{1}{2}$.}
}
